\section{Animation}\label{animation}

\subsection{Flight animation - workflow}\label{flight-animation}

This section explains the necessary steps required to create flight animation. \emph{Flight animation} in Mandelbulber is like a camera motion track recorded in some kind of flight simulator. The camera can go through interiors of fractal objects.
Recording is normally initially done by navigating with the image set at a low resolution (e.g., 320 x 240). After previewing and making any changes, the final rendering is undertaken with the image resolution set to a suitable higher size.

The parameters of every single animation frame recorded in Flight Animation mode can be edited in the animation table.

\underline{Workflow:}

\begin{enumerate}
	\item Define fractal object (or many objects). 
	
	Create a fractal with interesting features for the flight animation, (e.g., interesting shapes, geometric structure, texture, coloring, possibly holes where the camera can navigate into).
	
	It is advisable to select a fractal object that is relatively fast to render. Using a fast-rendering fractal at a low resolution results in the navigation and flight path recording happening at almost real time (or slow motion). A fast-rendering fractal will also increase the speed of the final frame rendering process.
	
	\item Place the camera at the point where the flight is to start from.
			
	\item Set a low image resolution. At a low image resolution, the frames-per-second value can be set higher for the flight path recording. It is reasonable to use a resolution like 320x240 or 160x120.
	
	\item Disable all effects which can slow down rendering, like ambient occlusion, reflections, transparency, volumetric lights, etc. All these effects can be re-enabled before commencing the final rendering of the animation.
	
	\item Open \emph{Flight animation} editor. It can be opened from top pull-down-menu by activating \emph{View} / \emph{Show animation dock}. The dock will appear at the bottom of the application window, with \emph{Flight animation (every frame)} tab (showed on picture \ref{flight-animation-dock}).
	
	\simpleImageWithCaption75Width{img/manual/media/flight_animation_dock.png}
	{Flight animation dock}
	{flight-animation-dock}{h}
	
	\item Set parameters of animation
	
	\begin{description}
		\item[speed] defines how fast the camera will fly. This parameter can be changed during the recording of the flight path.
		\item[inertia] defines how heavy is the camera. A heavy camera will make a smoother motion, but it will be more difficult to control.
		\item[rotation speed] defines the reaction  speed of mouse pointer movements. A higher value will allow for the camera to be turned faster.
		\item[roll speed] defines the speed of the reaction for the Z and X keys, which rotate the camera.
		\item[speed control] defines how the speed of the camera will be controlled.
		\begin{itemize} 
			\item In \emph{Relative to distance} mode, the camera speed will decrease relatively to the nearness of the fractal surface. This mode will help you to not collide with the fractal. In this mode, you can control the relative speed by the speed parameter and by the mouse buttons.
			\item In \emph{Constant mode}, the camera speed is only controlled by the speed parameter and the mouse buttons.
		\end{itemize}
		\item[second per frame] defines the frame rate during flight path recording. Higher values give slower rendering but images are more detailed. The value of this parameter is used only during recording, and is ignored during the rendering of the final animation.
		\item[path for images] defines where the rendered animation frames will be stored.
		\item[image type] defines the image format for saving the rendered frames. Detailed settings for image format are in \emph{File} / \emph{Program preferences}.
		\item[show thumbnails] enables previews of frames in animation table.
		\item[add flight and rotation speed to parameters] enables possibility to continue recording of animation after recording is completely stopped. 
	\end{description}

	\item Press \emph{Record flight path} button. After 3 seconds, recording will be started. During this 3-seconds-waiting time the mouse cursor should be placed in the center of image.
	
	\item Use mouse pointer movements to turn camera left / right and up / down. Camera behaves like airplane in flight simulator.
	
	Use Z and X key to rotate the camera.
	
	Use arrow keys to move camera left / right / up / down. Without \emph{Shift} key, the camera still goes forward (movement at 45-degree angle). With \emph{Shift} key, the camera is not moving forward (movement at straight angle).
	
	Use the left mouse button to increase flight speed, or the right button to decrease speed.
	
	\item Press \emph{space} key to pause recording. When recording is paused, animation parameters can be changed. 
	
	\item Press \emph{STOP} button to stop recording. It is good to pause the recording using the (space key) before stopping. This is because moving the mouse pointer towards the \emph{STOP} button can turn the camera.
	
	\item Recording of animation can be continued if \emph{add flight and rotation speed to parameters} is enabled. If \emph{Continue recording} is pressed, the recording of flight will resume from the point stored in the last frame. The camera linear and rotation speed will be maintained.
	
	\item Increase image resolution and enable all required effects. There can be added light sources, fog, materials, textures, etc.
	
	\item Press \emph{Render flight animation} to commence final rendering  process for the animation. This can take a very long time depending on the image resolution and the number of frames.
	
	Rendering of the animation can be stopped at any time and continued later. When \emph{Render flight animation} is pressed, there will be rendered only the frames which are missing from the image frame folder. Any existing rendered frames will be skipped.
	
\end{enumerate}
	
	\subsection{Flight animation - more options}
	
	\subsubsection{Adding more parameters to animation}
	
	It is possible to animate almost any parameter in Mandelbulber. Each edit field has an assigned parameter. If you place the mouse pointer on an edit field, a tooltip will be displayed (figure \ref{example-tooltip-with-parameter-name}).
	
	\simpleImageWithCaption75Width{img/manual/media/tooltip_with_parameter_name.png}
	{Example tooltip with parameter name}
	{example-tooltip-with-parameter-name}{h}
		
	Below the parameter description there is the \emph{Parameter name} (in this example it is \emph{glow\_intensity}).
	
	The parameter can be added to \emph{Flight animation}. Right click on an edit field, then use \emph{Add to flight animation} from context menu, and the parameter will appear in the animation table  (picture \ref{flight-animation-table}).
			
	\subsubsection{Editing animation in the table}
	
	It is possible to edit every single animation frame in the animation table. Every cell in the table is editable.  When a value is being edited, the preview is refreshed.
	
		\simpleImageWithCaption75Width{img/manual/media/editing_parameters_in_flight_table.png}
	{Animation table with all recorded frames and one additional parameter}
	{flight-animation-table}{h}
	
	Picture \ref{flight-animation-table} shows changes to \emph{main\_glow\_intensity}. On frame 2 it is changed to 10 and on frame 5 it is changed to 50.
	
	To get a smooth change of value within the selected range of frames, there is an option to do linear interpolation of values. 
	
	\underline{Example:} We would like to have smooth transition of glow intensity starting from frame 30 (\emph{glow} = 0.2) and ending on frame 100 (\emph{glow} = 10).
	\begin{itemize}
		\item Right click in the table on cell for glow intensity at frame 100
		\item Use option  \emph{Interpolate next frames}
		\item \emph{Last frame number} set to 100
		\item \emph{Value for last frame} set to 10
		\item When you press OK, you should see that in the table the glow intensity parameter is increasing to 10 from frame 30 to frame 100
	\end{itemize} 
	
	To cut the animation at the start or at the end, there is an option to delete a range of frames. If you right click on the frame number, there are two options:
	\begin{itemize}
		\item \emph{Delete all frames to here} - deletes all frames from first to selected frame
		\item \emph{Delete all frames from here} - deletes all frames from selected frame to the end
	\end{itemize}
	
	
	%TODO adding of parameters to animation, manual editing, deleting of frames
	


\subsection{Keyframe animation - workflow}\label{keyframe-animation}

\emph{Keyframe animation} works by defining a set of keyframes, each containing a complete snapshot of all animated parameters (camera position, fractal settings, light positions, material properties, etc.). Mandelbulber then interpolates between these keyframes to generate every intermediate frame.

This approach is ideal for animations that require precise control over the scene at specific points, such as camera movements, fractal morphing, or lighting transitions.

\subsubsection{Opening the Keyframe animation editor}

Open the \emph{Animation} dock from the top pull-down menu by activating \emph{View} / \emph{Show animation dock}. Switch to the \emph{Keyframe animation (only keyframes)} tab (picture \ref{keyframe-animation-dock}).

\simpleImageWithCaption75Width{img/manual/media/keyframe_animation_dock.png}
{Keyframe animation dock}
{keyframe-animation-dock}{H}

The dock is divided into three main areas: the control buttons on the left, the animation parameters panel below the buttons, and the keyframe table with the value chart on the right.

\subsubsection{Adding parameters to animation}

Before creating keyframes, you must tell Mandelbulber which parameters you want to animate. By default, no parameters are tracked.

To add a parameter, right-click on its edit field in any panel (fractal, camera, lights, materials, etc.) and select \emph{Add to keyframe animation} from the context menu. The parameter will appear as a new row in the keyframe table.

To add all currently non-default parameters at once, click the \emph{Add all parameters} button in the animation parameters panel.

\subsubsection{Adding keyframes}

\begin{enumerate}
	\item Set up your scene (fractal, camera, lights, materials, etc.).
	
	\item Add the parameters you want to animate (see above).
	
	\item Click the \emph{Add keyframe} button to append a keyframe at the end of the animation. The current values of all tracked parameters are captured.
	
	\item Modify the parameters in the main window (e.g., move the camera, change the fractal power, adjust light positions).
	
	\item Click the \emph{Modify keyframe} button to update the selected keyframe with the new values.
	
	\item Repeat steps 2--4 to build your animation.
	
	\item To insert a keyframe at a specific position, select the target column in the table and click \emph{Insert keyframe}. The new keyframe will be placed before the selected one.
	
	\item The \emph{Randomize} button opens the Settings Randomizer dialog for the currently selected keyframe. It allows you to randomly modify the tracked parameter values with three strength levels (\emph{Slight}, \emph{Medium}, \emph{Heavy}). The dialog includes a live preview render, a list of parameters to keep unchanged, and options for saving/loading randomizer presets. After confirming, the randomized values are applied to the selected keyframe.
	
	\item The \emph{Add all parameters} button adds all currently non-default parameters to the animation table at once, saving the need to add them individually via the context menu.
\end{enumerate}

Keyboard shortcuts: press \emph{K} to add a keyframe, \emph{M} to modify the selected keyframe.

\subsubsection{Animation parameters panel}

The parameters panel (left side of the dock) controls the overall animation timing and output (picture \ref{keyframe-animation-dock}):

\begin{description}
	\item[frames per keyframe] defines how many interpolated frames are generated between each keyframe. The total number of animation frames is calculated as: $(\text{number of keyframes} - 1) \times \text{frames per keyframe} + 1$. The last keyframe always has exactly one sub-frame (it is the endpoint). Default: 100.
	
	\item[frames per second] sets the animation playback speed. Default: 30 fps.
	
	\item[path for images] defines the output directory for rendered frames. Click \emph{Select folder} to choose a location.
	
	\item[image type] selects the image format for rendered frames (PNG, JPG, EXR, TIFF). Detailed format settings are in \emph{File} / \emph{Program preferences}.
	
	\item[first / last frame to render] limits the range of frames to render. Set to 0 to render all frames.
	
	\item[loop] enables seamless looping of the animation by interpolating between the last and first keyframes. Requires at least 3 keyframes.
	
	\item[show keyframe preview thumbnails] enables small preview images of each keyframe in the animation table.
	
	\item[show camera / target / light path] toggles visualization of the camera, target, and light movement paths in the 3D viewport.
\end{description}

\subsubsection{The keyframe table}

The keyframe table (picture \ref{keyframe-table}) displays all keyframes as columns and all animated parameters as rows. Each cell contains the parameter value for that specific keyframe.

\simpleImageWithCaption75Width{img/manual/media/keyframe_table_with_interpolation.png}
{Keyframe animation table}
{keyframe-table}{H}

The table has reserved rows at the top (parameter names) and reserved columns at the left (frame numbers). The background color of each cell indicates the interpolation type for that parameter (see below).

Every cell in the table is directly editable. When a value is being edited, the preview is refreshed in real time.

\subsubsection{Interpolation types}

Each animated parameter row has its own interpolation type, which controls how values transition between keyframes. To set the interpolation type, right-click on the parameter name in the row header (picture \ref{keyframe-table}) and choose from the context menu.

The interpolation type is indicated by the background color of cells in that row:

\begin{description}
	\item[None] (white) -- no interpolation; the value remains constant until the next keyframe.
	\item[Linear] (gray) -- straight-line interpolation between keyframe values.
	\item[Linear angle] (gray/blue) -- linear interpolation with angle-aware rotation, designed for quaternion parameters.
	\item[CatMulRom] (green) -- Catmull-Rom spline: produces a smooth C1 continuous curve passing through all keyframe values.
	\item[CatMulRom angle] (green/blue) -- Catmull-Rom spline with angle-aware rotation.
	\item[Akima] (red) -- Akima spline: smooth and monotonic, avoids overshoot. Useful for parameters that should not oscillate.
	\item[Akima angle] (red/blue) -- Akima spline with angle-aware rotation.
	\item[Cubic] (yellow) -- Cubic spline: produces a smooth C2 continuous curve.
	\item[Cubic angle] (yellow/blue) -- Cubic spline with angle-aware rotation.
	\item[Steffen] (blue) -- Steffen monotonic spline: prevents overshoot while maintaining smoothness.
	\item[Steffen angle] (blue/blue) -- Steffen spline with angle-aware rotation.
\end{description}

The \emph{angle} variants are specifically designed for quaternion-based parameters (camera rotation) to avoid issues with short-path interpolation.

\subsubsection{Editing values in the table}

\begin{description}
	\item[Direct editing] Click any cell to edit its value. The preview updates as you type.
	
	\item[Copy value to all keyframes] Right-click a cell and select \emph{Copy value to all keyframes} to copy that value to the same parameter across all keyframes.
	
	\item[Multiply / Increase values by] Select multiple cells (drag or shift-click), then right-click and choose \emph{Multiply values by...} or \emph{Increase values by...} to apply a scalar operation to all selected cells.
	
	\item[Interpolate keyframes] Right-click a cell and select \emph{Interpolate next keyframes} to create a smooth transition to a target value across subsequent keyframes.
	
	\item[Render this keyframe] Right-click a keyframe column header and select \emph{Render this keyframe} to render a preview of that frame.
	
	\item[Delete keyframes] Right-click a keyframe column header to delete that keyframe, or to \emph{Delete all keyframes to here} / \emph{Delete all keyframes from here} for bulk deletion.
	
	\item[Insert keyframe in between] Right-click a keyframe column header (not the first or second) and select \emph{Insert keyframe in between} to add a new keyframe at that position.
	
	\item[Remove parameter from animation] Right-click a parameter name row header and select \emph{Remove 'parameter\_name' from animation}.
	
	\item[Refresh all thumbnails] Right-click the reserved header row at the top and select \emph{Refresh all thumbnails} to regenerate preview images.
\end{description}

\subsubsection{Value chart}

Above the keyframe table is the value chart (picture \ref{keyframe-chart}), which visualizes the values of animated parameters across all frames. This helps identify the shape of parameter transitions and detect potential issues such as sudden jumps.

\simpleImageWithCaption75Width{img/manual/media/keyframe_value_chart.png}
{Value chart}
{keyframe-chart}{H}

Use the \emph{+} and \emph{−} buttons to zoom in and out. The frame slider below the chart allows navigation through interpolated frames, and the preview updates accordingly.

\subsubsection{Validation and cleanup}

The \emph{Validation and cleaning up} panel provides tools to ensure animation quality:

\begin{description}
	\item[Validate before render] enables automatic collision checking before rendering. If a collision is detected, the animation will not start.
	
	\item[Validate] (button) manually checks all interpolated frames for collisions between the camera and the fractal. A collision occurs when the distance from the camera to the fractal surface is lower than the \emph{collision distance} threshold.
	
	\item[Collision distance] sets the minimum allowed distance between the camera and the fractal surface. Reduce this value to allow closer approaches.
	
	\item[Set the same camera target distance for all keyframes] (button) equalizes the distance between the camera and its target across all keyframes. This keeps the camera rotation smooth throughout the animation without modifying the camera position or movement path. The target distance value can be set in the adjacent field.
\end{description}

\subsubsection{Rendering the animation}

\begin{enumerate}
	\item Verify the animation using the value chart and frame slider.
	
	\item Optionally run \emph{Validate} to check for collisions.
	
	\item Set the output path and image format in the animation parameters panel.
	
	\item Click \emph{Render animation}. The program will render only frames that are missing from the output folder. Existing frames are skipped, allowing interrupted renders to be resumed.
	
	\item When rendering is complete, click \emph{Show Animation} to play back the rendered frames.
\end{enumerate}

\subsubsection{Exporting between animation modes}

Keyframe animations can be converted to Flight animation and vice versa:

\begin{description}
	\item[Export to Flight] (button) exports all interpolated keyframe frames to the \emph{Flight animation (every frame)} tab. Each keyframe is expanded into individual per-frame flight data, allowing fine-grained editing of every frame in the Flight animation editor (or with an external CSV editor).
	
	\item[Export to Keyframes] (button, on the Flight animation tab) converts a recorded flight path into keyframes. The \emph{frames per keyframe} parameter on the Keyframe animation tab determines how many flight frames are sampled to create new keyframes.
\end{description}

\subsubsection{Animation storage format}

Keyframe animation data is stored in the Mandelbulber settings file (\emph{.fract}) in CSV format within the \texttt{[keyframes]} section. This means animations can be edited with external tools such as spreadsheet editors or text editors.

The first row of the \texttt{[keyframes]} section contains the parameter names. Each subsequent row represents one keyframe with its parameter values. The last row specifies the interpolation type for each parameter (e.g., \emph{morphAkima}, \emph{morphLinear}).
